% GENERATED by scripts/ws_i_prompts_appendix.py -- do not edit by hand.
% Prompts are imported from the harness, so they cannot drift from the code.
\section{Exact prompts}\label{app:prompts}
Every prompt below is rendered by importing the harness's own builder and filling it with one small worked example, so what appears here is what the code sends. Module text is abbreviated for space; nothing else is edited.

\paragraph{RepairBench no-spec arm (Table~\ref{tab:recovery}).}
The diagnostic is the verbatim tool output that defined the bug. This is a one-shot repair prompt; the returned module is post-scored by the class-appropriate gate.
\textit{System.}
\begin{Verbatim}[fontsize=\scriptsize,breaklines,breakanywhere]
You are an expert hardware design assistant. The user reports a tool error on a SystemVerilog module. First write ONE short paragraph explaining the cause of the error in plain English. Then output the complete corrected SystemVerilog inside a single fenced code block tagged ```systemverilog. Do not add commentary after the code block. Preserve the original module name and port list unless the error specifically requires changing them.
\end{Verbatim}
\textit{User.}
\begin{Verbatim}[fontsize=\scriptsize,breaklines,breakanywhere]
A verification tool reported the following problem with this SystemVerilog module:

```
%Warning-WIDTH: TopModule.sv:4: Operator ASSIGNDLY expects 1 bit on the Assign RHS.
```

The failing module is:

```systemverilog
module TopModule(input clk, input rst_n, input a, output reg z);
  always @(posedge clk)
    if (rst_n) z <= 1'b0;
    else z <= a;
endmodule
```

Explain the cause and return the fixed module.
\end{Verbatim}


\paragraph{RepairBench specification arm (Table~\ref{tab:recovery}).}
Same system message and same builder as the arm above; only these lines are added. RepairBench's builder words things differently from the decomposition harness, which is why the two are shown separately.
\textit{User.}
\begin{Verbatim}[fontsize=\scriptsize,breaklines,breakanywhere]
The module is supposed to implement this specification:
z is registered a, cleared while the active-low reset rst_n is asserted.
\end{Verbatim}


\paragraph{Decomposition arms, frozen primary (Table~\ref{tab:primary}).}
Arm $\mathrm{spec}0\,\mathrm{loc}0$, the no-signal baseline. The three remaining arms add the blocks below, unchanged otherwise.
\textit{System.}
\begin{Verbatim}[fontsize=\scriptsize,breaklines,breakanywhere]
You are an expert RTL debugging engineer. Repair the supplied
compile-clean SystemVerilog module. Return the complete replacement module,
from `module` through `endmodule`; never return a patch, explanation, testbench,
or placeholder. Preserve the given module interface and emit synthesizable RTL.
\end{Verbatim}
\textit{User.}
\begin{Verbatim}[fontsize=\scriptsize,breaklines,breakanywhere]
Repair the following compile-clean SystemVerilog module, which fails functional verification.

Buggy module:
```systemverilog
module TopModule(input clk, input rst_n, input a, output reg z);
  always @(posedge clk)
    if (rst_n) z <= 1'b0;
    else z <= a;
endmodule
```

Return only the complete corrected `TopModule`, including its unchanged interface.
\end{Verbatim}


\paragraph{Decomposition arms, post hoc capability curve (Figure~\ref{fig:curve}).}
The four user messages are byte-identical to the frozen-primary messages above, which this script asserts on every run. Only the system message differs, because the historical curve runner renders it as a single line.
\textit{System.}
\begin{Verbatim}[fontsize=\scriptsize,breaklines,breakanywhere]
You are an expert RTL debugging engineer. Repair the supplied compile-clean SystemVerilog module. Return the complete replacement module, from `module` through `endmodule`; never return a patch, explanation, testbench, or placeholder. Preserve the given module interface and emit synthesizable RTL.
\end{Verbatim}


\paragraph{Shared What (specification) block, added for $\mathrm{spec}1$.}

\textit{User.}
\begin{Verbatim}[fontsize=\scriptsize,breaklines,breakanywhere]
Behavioral specification:
z is registered a, cleared while the active-low reset rst_n is asserted.
\end{Verbatim}


\paragraph{Shared Where (oracle localization) block, added for $\mathrm{loc}1$.}
The line number and the original broken line only; never the fix.
\textit{User.}
\begin{Verbatim}[fontsize=\scriptsize,breaklines,breakanywhere]
Oracle localization (this identifies the buggy source line but does not give the fix):
Line 3: `if (rst_n) z <= 1'b0;`
\end{Verbatim}


\paragraph{Counterexample revision pair (Table~\ref{tab:primary}, third contrast).}
The concrete arm. The generic arm is byte-identical except that the trace block is replaced by the single sentence ``The previous candidate still did not pass verification.'', which is what isolates the counterexample from re-exposure.
\textit{System.}
\begin{Verbatim}[fontsize=\scriptsize,breaklines,breakanywhere]
You are an expert RTL debugging engineer. Repair the supplied
compile-clean SystemVerilog module. Return the complete replacement module,
from `module` through `endmodule`; never return a patch, explanation, testbench,
or placeholder. Preserve the given module interface and emit synthesizable RTL.
\end{Verbatim}
\textit{User.}
\begin{Verbatim}[fontsize=\scriptsize,breaklines,breakanywhere]
Revise a previous repair attempt for this original problem.

Original buggy module:
```systemverilog
module TopModule(input clk, input rst_n, input a, output reg z);
  always @(posedge clk)
    if (rst_n) z <= 1'b0;
    else z <= a;
endmodule
```

Previous candidate:
```systemverilog
module TopModule(input clk, input rst_n, input a, output reg z);
  always @(posedge clk)
    if (!rst_n) z <= 1'b0;
    else z <= 1'b1;
endmodule
```

The verifier produced this concrete trace from initialization through the first divergence. Values include the driven inputs, reference expected outputs, and candidate outputs:
```text
cycle 0: rst_n=0 a=1 | expected z=0 got z=1
cycle 1: rst_n=1 a=1 | expected z=1 got z=1
```

Return only the complete corrected `TopModule`.
\end{Verbatim}


\paragraph{Mined-failure arms (Table~\ref{tab:multimodel}).}
This is \code{nospec}. The \code{spec} arm is the same message with the specification block appended, and \code{speconly} is the same message with the broken RTL omitted and the specification kept, so it carries no module to edit and no divergence trace.
\textit{System.}
\begin{Verbatim}[fontsize=\scriptsize,breaklines,breakanywhere]
You repair broken SystemVerilog.
\end{Verbatim}
\textit{User.}
\begin{Verbatim}[fontsize=\scriptsize,breaklines,breakanywhere]
Broken RTL (`TopModule.sv`):
```systemverilog
module TopModule(input clk, input rst_n, input a, output reg z);
  always @(posedge clk)
    if (rst_n) z <= 1'b0;
    else z <= a;
endmodule
```

The module compiles but is functionally incorrect. Return the complete corrected module (module ... endmodule), interface and module name unchanged.
\end{Verbatim}


\paragraph{Redaction control (\S\ref{sec:c2realbug}).}
Applied to the specification only. The redacted specifications are not released; only aggregate summaries are included in this artifact.
\textit{System.}
\begin{Verbatim}[fontsize=\scriptsize,breaklines,breakanywhere]
You are redacting a hardware module specification for a controlled experiment. Rewrite the specification so it still states the module's PURPOSE and its INPUT/OUTPUT INTERFACE, but REMOVES or GENERALIZES every detail that reveals the exact correct logic: specific boolean operators, comparisons, equalities, polarities (active-high/low), literal constants and reset values, exact state-transition tables, truth-table rows, and any phrase that names the precise function (e.g. replace 'a NOR gate' or 'z = (x^y) & x' with 'the specified combinational function', replace a full state table with 'the state machine described by the following interface'). Keep it a plausible, well-formed spec a designer could receive, just without the giveaway details. Do NOT add new information or the answer. Return ONLY the redacted specification text, no preamble.
\end{Verbatim}
\textit{User.}
\begin{Verbatim}[fontsize=\scriptsize,breaklines,breakanywhere]
(the specification to be redacted)
\end{Verbatim}

